\section{Varianz \& Standardabweichung (Kap.\ 3.3)}

\subsection{Funktionen von Zufallsvariablen}
\begin{merke}[{{Idee}}]
\begin{itemize}
\pfeil $X:\Omega\to\mathbb R^n$ diskret, $f:\Omega_X\to\mathbb R^m$ $\Rightarrow$ $Y=f(X)=f\circ X$ ist wieder eine ZV
\pfeil Verteilung von $X$ genügt, $(\Omega,\Sigma,P)$ nicht nötig
\end{itemize}
\end{merke}
\begin{formel}
$\displaystyle E[f(X)]=\sum_{x\in\Omega_X} f(x)\cdot p_X(x)$ \quad (Thm.\ 3.3.2)
\end{formel}
\begin{itemize}
\pfeil Nur anwendbar, falls $E[f(X)]$ existiert (abs.\ Konvergenz)
\pfeil Verteilung von $f(X)$ \emph{nicht} nötig: direkt $f(x)$ mit $p_X(x)$ gewichten
\pfeil Linearität: $E[a_1X_1+\dots+a_nX_n]=\sum a_iE[X_i]$ (auch ohne Unabhängigkeit)
\end{itemize}
\begin{bsp}[{{$X$ = Augenzahl Würfel, $f(x)=x^2$}}]
$E[X^2]=\tfrac16(1+4+9+16+25+36)=\tfrac{91}{6}\approx15.17$\\
Achtung: $E[X^2]\neq(E[X])^2=3.5^2=12.25$
\end{bsp}

\subsection{Definition}
\begin{formel}
$\mu=E[X]$\\[1pt]
$\displaystyle \operatorname{var}(X)=E\big[(X-\mu)^2\big]=\sum_{x\in\Omega_X}(x-\mu)^2\,p_X(x)\ \ge0$\\[1pt]
$\sigma(X)=\sqrt{\operatorname{var}(X)}\ \ge0$
\end{formel}
\begin{itemize}
\pfeil var = gewichtete Summe der quadrierten Abweichungen vom Erwartungswert
\pfeil $\sigma$ = typische Abweichung vom $\mu$ (gleiche Einheit wie $X$)
\pfeil Je grösser var, desto breiter streut $p_X$ um $\mu$
\pfeil Existenz: $E[X^2]$ existiert $\Leftrightarrow$ var$(X)$ existiert
\end{itemize}

\subsection{Rechenregeln}
\begin{formel}
\textbf{Verschiebungssatz:}\quad
$\operatorname{var}(X)=E[X^2]-(E[X])^2$
\end{formel}
\begin{formel}
$\operatorname{var}(aX+b)=a^2\operatorname{var}(X)$\qquad
$\sigma(aX+b)=|a|\,\sigma(X)$
\end{formel}
\begin{itemize}
\pfeil Verschiebung $+b$: Streuung \emph{unverändert}
\pfeil Skalierung $\cdot a$: $\sigma$ mal $|a|$, var mal $a^2$
\pfeil $\operatorname{var}(X)\ge0$, also $E[X^2]\ge(E[X])^2$
\pfeil $\operatorname{var}(X)=0\Leftrightarrow X$ konstant
\end{itemize}


\begin{merke}[{{Rechenschema Varianz (Tabelle)}}]
\begin{enumerate}[label=\arabic*.,leftmargin=*]
\item Verteilung $p_X(x)$ aufstellen (Summe $=1$ prüfen)
\item $\mu=E[X]=\sum x\,p_X(x)$
\item $E[X^2]=\sum x^2\,p_X(x)$
\item $\operatorname{var}=E[X^2]-\mu^2$, \ $\sigma=\sqrt{\operatorname{var}}$
\end{enumerate}
\begin{itemize}
\pfeil Alternativ: $\sum (x-\mu)^2p_X(x)$ (rechnet sich mühsamer)
\end{itemize}
\end{merke}

\begin{bsp}[{{Scrabble-Stein, Punkte $X$ (102 Steine)}}]
\begin{tabular}{@{}c|cccccccc@{}}
$x$&0&1&2&3&4&6&8&10\\\midrule
$p_X$&$\frac{2}{102}$&$\frac{64}{102}$&$\frac{13}{102}$&$\frac{8}{102}$&$\frac{7}{102}$&$\frac{4}{102}$&$\frac{2}{102}$&$\frac{2}{102}$
\end{tabular}\\[2pt]
$E[X]=\frac{64+26+24+28+24+16+20}{102}=\frac{101}{51}\approx1.98$\\
$E[X^2]=\frac{64+52+72+112+144+128+200}{102}=\frac{386}{51}\approx7.57$\\
$\operatorname{var}=\frac{386}{51}-\big(\frac{101}{51}\big)^2\approx3.65$,\ \ $\sigma\approx1.91$
\end{bsp}

\begin{bsp}[{{Würfel: $X$ = Augenzahl}}]
$\mu=3.5$,\ $E[X^2]=\frac{91}{6}$\\
$\operatorname{var}=\frac{91}{6}-\frac{49}{4}=\frac{35}{12}\approx2.92$,\ $\sigma\approx1.71$
\end{bsp}

\begin{bsp}[{{Summe zweier Würfel $S=X_1+X_2$}}]
$\mu=E[S]=7$\\
$\operatorname{var}(S)=\sum(s-7)^2p_S(s)=\frac{2\,(25\cdot1+16\cdot2+9\cdot3+4\cdot4+1\cdot5)}{36}+0$\\
$=\frac{2\cdot105}{36}=\frac{35}{6}\approx5.83$,\ $\sigma\approx2.42$\\
$\Rightarrow$ var$(X_1)+$var$(X_2)=\frac{35}{12}\cdot2$ (gilt für \emph{unabhängige} ZV)
\end{bsp}

\subsection{Chebyshev-Ungleichung}
\begin{formel}
$\displaystyle P\big(|X-\mu|\ge d\big)\le\frac{\operatorname{var}(X)}{d^2}\qquad(d>0)$
\end{formel}
\begin{formel}
$d=k\sigma$:\quad $\displaystyle P\big(|X-\mu|\ge k\sigma\big)\le\frac1{k^2}$
\end{formel}
\begin{itemize}
\pfeil $k=2$: $P(|X-\mu|\ge2\sigma)\le\frac14$\ \ ($k=3$: $\le\frac19$)
\pfeil Gilt für \emph{jede} Verteilung (nur $\mu$, var nötig)
\pfeil Obere Schranke, oft grob
\pfeil Komplement: $P(|X-\mu|<d)\ge1-\frac{\operatorname{var}}{d^2}$
\pfeil Nur sinnvoll für $d>\sigma$ (sonst Schranke $\ge1$)
\end{itemize}
\begin{bsp}[{{Chebyshev anwenden}}]
$\mu=10$, $\sigma=2$, gesucht $P(|X-10|\ge6)$\\
$d=6=3\sigma\Rightarrow P\le\frac{4}{36}=\frac19$
\end{bsp}

\subsection{Geometrische Verteilung}
\begin{merke}[{{$X\sim\mathrm{Geom}(p)$, $0<p<1$}}]
\begin{itemize}
\pfeil $X$ = Anzahl Versuche bis 1.\ Erfolg, $n\in\mathbb N$
\pfeil $p_X(n)=(1-p)^{n-1}p$
\end{itemize}
\end{merke}
\begin{formel}
$\displaystyle E[X]=\frac1p\qquad
\operatorname{var}(X)=\frac{1-p}{p^2}\qquad
\sigma(X)=\frac{\sqrt{1-p}}{p}$
\end{formel}
\begin{formel}
$\displaystyle P(X\le k)=1-(1-p)^k$
\end{formel}
\begin{itemize}
\pfeil Je kleiner $p$, desto grösser $\mu$ \emph{und} Streuung
\end{itemize}
\begin{bsp}[{{Aufg.\ 3.3.10: Gewinnwahrsch.\ $p=\frac1{10}$}}]
$E[X]=\frac1p=10$ Spiele\\
$\operatorname{var}=\frac{0.9}{0.01}=90$,\ \ $\sigma=\sqrt{90}\approx9.49$\\
$\Rightarrow$ $\sigma\approx\mu$: sehr breite Streuung
\end{bsp}

\subsection{Prüfungsbeispiele (Muster)}
\begin{bsp}[{{Roulette-Gewinn $X\in\{-1,+1\}$, $p=\frac{20}{38},\frac{18}{38}$}}]
$\mu=-\frac{2}{38}\approx-0.053$,\ $E[X^2]=1$\\
$\operatorname{var}=1-\mu^2\approx0.997$,\ $\sigma\approx1$
\end{bsp}
\begin{bsp}[{{Lineare Transformation}}]
$\sigma(X)=3$, $Y=-2X+5$\\
$\operatorname{var}(Y)=(-2)^2\cdot9=36$,\ $\sigma(Y)=2\cdot3=6$\\
(die $+5$ ändert nichts)
\end{bsp}
\begin{merke}[{{Checkliste}}]
\begin{itemize}
\pfeil $\sum p_X=1$?\ \ $\mu$ zuerst, dann $E[X^2]$
\pfeil $E[X^2]\ne(E[X])^2$
\pfeil var $\ge0$ als Kontrolle (negativ $\Rightarrow$ Rechenfehler)
\pfeil Faktor $a$: $a^2$ bei var, $|a|$ bei $\sigma$
\pfeil Verteilung bekannt (Geom)? $\to$ Formel statt Rechnen
\end{itemize}
\end{merke}
